% problem-type: laplace-vanh-khan
% order: 80
% Nguon: De thi MAT2306 (PTDHR 1), HK2 2024-2025, K67, Cau 3. (co dap an chinh thuc)
% Bo cuc: De vi du -> Cong thuc -> De + Bai giai

\section{Laplace trong (góc phần tư) vành khăn --- phản xạ Schwarz}

% ---------------------------------------------------------------
\subsection*{Đề bài ví dụ}
% ---------------------------------------------------------------
Xét bài toán biên cho phương trình Laplace trong góc phần tư vành khăn
\[
  u_{xx}+u_{yy}=0,\quad x>0,\ y>0,\ 1<x^2+y^2<4,
\]
với điều kiện biên $u_x(0,y)=u_y(x,0)=0$ khi $1<x,y<2$, và
\[
  u(x,y)=2x\ \text{khi } x^2+y^2=1,\qquad u(x,y)=y\ \text{khi } x^2+y^2=4.
\]
(a) Dùng phản xạ Schwarz qua trục $x=0$ rồi qua trục $y=0$, bài toán chuyển thành bài toán nào
trong vành khăn? Từ đó chứng minh $0<u<2$. \ (b) Giải bài toán.

\emph{\small(Nguồn: Đề thi MAT2306 --- PTĐHR 1, HK2 2024--2025, K67, Câu 3.)}

% ---------------------------------------------------------------
\subsection*{Nhận ra dạng này khi nào?}
% ---------------------------------------------------------------
\begin{itemize}
  \item Miền là \textbf{vành khăn} (không chứa gốc) $\Rightarrow$ nghiệm tách biến giữ CẢ $r^{n}$ lẫn $r^{-n}$, và $\ln r$.
  \item Cho \textbf{góc phần tư} vành khăn với Neumann ($u_x$ hoặc $u_y=0$) trên các cạnh thẳng
  $\Rightarrow$ \textbf{phản xạ Schwarz} (chẵn) để gập thành vành khăn đủ rồi tách biến.
\end{itemize}
\begin{meobox}
Cạnh thẳng Neumann $\Rightarrow$ thác triển \emph{chẵn}; mỗi lần phản xạ qua một trục nhân đôi góc.
Hai cạnh $\theta=0,\tfrac{\pi}{2}$ với $v_\theta=0$ $\Rightarrow$ chuỗi chỉ chứa $\cos(2n\theta)$.
\end{meobox}

% ---------------------------------------------------------------
\subsection*{Công thức \& phương pháp}
% ---------------------------------------------------------------
\subsubsection*{Phản xạ Schwarz (gập miền)}
% technique: phan-xa-schwarz
\textbf{Phản xạ Schwarz (gập miền điều hòa).}
Hàm điều hòa trên miền, có biên \emph{phẳng/cung} với:
\begin{itemize}
  \item Dirichlet $u=0$ $\Rightarrow$ thác triển \textbf{lẻ} qua biên đó (vẫn điều hòa).
  \item Neumann $\partial_\nu u=0$ $\Rightarrow$ thác triển \textbf{chẵn} qua biên đó.
\end{itemize}
Dùng để gập miền nhỏ thành miền đối xứng đầy đủ rồi giải. Ví dụ: góc phần tư vành khăn với
$u_x=0$ trên trục tung, $u_y=0$ trên trục hoành $\Rightarrow$ phản xạ chẵn hai lần thành \textbf{vành khăn đủ}.


\subsubsection*{Tách biến trong tọa độ cực}
% method: tach-bien-cuc
\begin{dieukienbox}
Dùng khi: $\Delta u=0$ (Laplace) hoặc $\Delta u=$ hằng (Poisson) trên \textbf{quạt / vành khăn / đĩa},
biên gồm cung tròn ($r=$ const) và cạnh thẳng ($\theta=$ const).
\end{dieukienbox}

\begin{nhobox}
\textbf{Laplacian cực:} $\Delta u=u_{rr}+\frac1r u_r+\frac{1}{r^2}u_{\theta\theta}$. Tách $u=R(r)\Theta(\theta)$:
\begin{itemize}
  \item Góc: $\Theta''+\mu^2\Theta=0\Rightarrow \Theta=\cos(\mu\theta)$ hoặc $\sin(\mu\theta)$ ($\mu$ từ biên ở 2 cạnh $\theta$).
  \item Bán kính (Euler): $r^2R''+rR'-\mu^2 R=0\Rightarrow R=r^{\mu},\,r^{-\mu}$ (hoặc $a+b\ln r$ khi $\mu=0$).
\end{itemize}
Nghiệm tổng quát:
$\displaystyle u=a_0+b_0\ln r+\sum_{n\ge1}\big(a_n r^{\mu_n}+b_n r^{-\mu_n}\big)\Theta_n(\theta).$
\end{nhobox}

\begin{meobox}
\begin{itemize}
  \item Cạnh Neumann $u_\theta=0$ tại $\theta=0$ $\Rightarrow$ chọn $\cos(\mu\theta)$; Dirichlet $u=0$ $\Rightarrow$ chọn $\sin(\mu\theta)$. Cạnh còn lại định $\mu_n$.
  \item \textbf{Đĩa} (chứa gốc): bỏ $r^{-\mu},\ln r$ (chặn tại $0$). \textbf{Vành khăn}: giữ cả hai.
  \item Poisson $\Delta u=1$: cộng nghiệm riêng $r^2/4$. Hệ số $a_n,b_n$ lấy từ khai triển Fourier dữ liệu trên cung.
\end{itemize}
\end{meobox}
\emph{(Laplacian theo tọa độ: Phụ lục B.)}

% ---------------------------------------------------------------
\subsection*{Đề bài \& bài giải}
% ---------------------------------------------------------------
\paragraph{Nhắc lại đề.} $\Delta u=0$ trong góc phần tư vành khăn $\{x,y>0,\ 1<x^2+y^2<4\}$;
$u_x(0,y)=u_y(x,0)=0$; $u=2x$ trên $r=1$, $u=y$ trên $r=2$.

\paragraph{(a) Phản xạ Schwarz $\to$ vành khăn + chặn $0<u<2$.}
Vì $u_x(0,y)=0$ ta thác triển \textbf{chẵn} theo $x$, rồi vì $u_y(x,0)=0$ thác triển chẵn theo $y$,
được hàm $U$ trên \textbf{toàn vành khăn} $1<x^2+y^2<4$ thỏa
\[
  U_{xx}+U_{yy}=0,\quad U=2|x|\ \text{trên } r=1,\quad U=|y|\ \text{trên } r=2.
\]
Trên biên: $0\le 2|x|\le2$ ($r=1$) và $0\le|y|\le2$ ($r=2$); hằng số là hàm điều hòa, nên theo
nguyên lý so sánh $0\le U\le2$. Dùng nguyên lý cực đại \emph{mạnh} ($U$ điều hòa, khác hằng trên biên)
$\Rightarrow 0<U<2$ trong miền.

\paragraph{(b) Bước 1 --- đổi sang cực, đọc biên.} Đặt $x=r\cos\theta,\ y=r\sin\theta$, $v(r,\theta)=U$.
Thác triển chẵn $\Rightarrow$
\[
  v_\theta(r,0)=v_\theta(r,\tfrac{\pi}{2})=0,\qquad v(1,\theta)=2\cos\theta,\quad v(2,\theta)=2\sin\theta.
\]

\paragraph{(b) Bước 2 --- nghiệm tổng quát (Neumann hai cạnh $\Rightarrow\cos(2n\theta)$).}
\[
  v(r,\theta)=a_0\ln r+b_0+\sum_{n\ge1}\big(a_n r^{2n}+b_n r^{-2n}\big)\cos(2n\theta).
\]

\paragraph{(b) Bước 3 --- khớp hai cung bằng Fourier.} Chiếu $2\cos\theta,\ 2\sin\theta$ lên hệ
$\{1,\cos(2n\theta)\}$ trên $[0,\tfrac{\pi}{2}]$:
\begin{align*}
  b_0&=\tfrac{2}{\pi}\!\int_0^{\pi/2}\!2\cos\theta\,d\theta=\tfrac{4}{\pi},\qquad
   a_0\ln2+b_0=\tfrac{2}{\pi}\!\int_0^{\pi/2}\!2\sin\theta\,d\theta=\tfrac{4}{\pi}\ \Rightarrow\ a_0=0;\\
  a_n+b_n&=\tfrac{4}{\pi}\!\int_0^{\pi/2}\!2\cos\theta\cos(2n\theta)\,d\theta=\frac{8(-1)^n}{\pi(1-4n^2)};\\
  4^n a_n+4^{-n}b_n&=\tfrac{4}{\pi}\!\int_0^{\pi/2}\!2\sin\theta\cos(2n\theta)\,d\theta=\frac{8}{\pi(1-4n^2)}.
\end{align*}
Giải hệ $2\times2$ cho mỗi $n$:
\[
  a_n=\frac{8\big(1-(-4)^{-n}\big)}{\pi(1-4n^2)(4^n-4^{-n})},\qquad
  b_n=\frac{8\big((-4)^{n}-1\big)}{\pi(1-4n^2)(4^n-4^{-n})}.
\]
\begin{nhobox}
\[
  \begin{aligned}
  v(r,\theta)=\frac{4}{\pi}+\sum_{n\ge1}\ &\frac{8}{\pi(1-4n^2)(4^n-4^{-n})}\\
  &\times\Big[r^{2n}\big(1-(-4)^{-n}\big)+r^{-2n}\big((-4)^n-1\big)\Big]\cos(2n\theta).
  \end{aligned}
\]
\end{nhobox}

% ---------------------------------------------------------------
\subsection*{Dạng bài lân cận}
% ---------------------------------------------------------------
\begin{itemize}
  \item Laplace trong \textbf{quạt/đĩa} (không phản xạ) --- xem mục riêng.
  \item \textbf{Poisson--Neumann} trong quạt ($\Delta u=$ hằng, kiểm tra tương thích) --- xem mục riêng.
\end{itemize}
